\cotefile{157-5}
\scanpage{293}
\mspage{209}

généralement, soit $J$ une \deleted{petite} catégorie telle que $\deleted{\unsure{c}}\operatorname{Ind}_{\pi_0} C$ soit stable par $\varprojlim$ de type $J$, et \deleted{soit} $c=\deleted{\operatorname{card}\unsure{\operatorname{Fl}|\operatorname{Ob}}}\inserted{\operatorname{card}\unsure{\operatorname{réd}|\operatorname{red}}}\,J$, supposons que $\pi^c=\pi$ (p.ex. $c$ fini). Alors $\operatorname{Filt}^{\pi}$ est stable \inserted{dans $M$} par $\varprojlim$ de \deleted{\unsure{cardinal}} type $J$. (Donc il en est de \unsure{même} des $M_{\pi}=\operatorname{Filt}^{\pi}$, si de plus $\pi^{\pi_0}=\pi$.) \deleted{\illegible{1}}

On utilise le critère $\unsure{a|A}''$ \inserted{de 4.7.1}, savoir $\operatorname{card}F(x)\leq\pi$ pour $\forall\,x\in C$, pour l'appartenance à $\operatorname{Filt}^{\pi}$. OPS $J$ réduite, avec \deleted{$\operatorname{card}J\leq$} $\operatorname{card}\operatorname{Ob}J\leq c$. Comme les \deleted{$\varprojlim$} \unsure{foncteurs} d'inclusion $\operatorname{Ind}_{\pi_0}(C)\longrightarrow C^\wedge$ commutent aux $\varprojlim$ quelconques, le calcul des $\varprojlim_{j\in J}x_j$ \deleted{\unsure{fibr}} dans $C^\wedge$ fibre par fibre montre que dans $C^\wedge$, \unsure{on a}

\[
X\overset{\deleted{\unsure{\mathrm{déf}}}}{=}\varprojlim_{j\in J}x_j\unsure{\subseteq|=}\underset{\substack{\rule{.3pt}{.5ex}\rule{4.8em}{.3pt}\rule{.3pt}{.5ex}\\ \text{produit dans }C^\wedge}}{\prod_{j\in\operatorname{Ob}J}x_j},
\]

\[
\text{d'où}\qquad \operatorname{card}X(x)\unsure{\leq|=}\prod_{j\in\operatorname{Ob}J}\underset{\substack{\rule{.3pt}{.5ex}\rule{4.5em}{.3pt}\rule{.3pt}{.5ex}\\ \leq\pi}}{\operatorname{card}x_j(x)}\leq\pi^c,
\]

d'où (sans autre hyp. sur $\pi$ que $\pi\geq\pi_{\unsure{1|0}}$)

\[
\deleted{\unsure{\varprojlim}}\varprojlim_{j\in J}x_j\in\operatorname{Ob}\operatorname{Filt}^{\pi^c},\tag*{(4.8.10.1)}
\]

d'où la conclusion voulue si $\pi^c=\pi$, \unsure{qed.|q.e.d.}
% archive sheet number: 292
\scanpage{294}
\mspage{210}

Nous allons prouver les implications supplémentaires

\[
b''\Longrightarrow b,\qquad c\Longrightarrow b,\qquad c''\overset{?}{\Longrightarrow}c,
\]

ce qui achèvera la démonstration.

L'implication $b''\Longrightarrow a$ est contenue dans 4.6.3. Pour l'implication $c\Longrightarrow b$, on note que la stabilité de \deleted{$\unsure{C}$} $M$ par petites $\unsure{\varinjlim|\underline{\lim}}$ équivaut à l'assertion que toute \inserted{petite} limite projective \deleted{$\unsure{\xi}$} de foncteurs $M\xrightarrow{\deleted{\unsure{F}}\unsure{\xi|\zeta}_i}\operatorname{Ens}$ représentables, est représentable. Or les objets de $M$ étant accessibles ($M$ étant accessible), les foncteurs représentables $\xi_i$ sont accessibles, donc aussi leur $\varprojlim\,\unsure{\xi|{}}$ (SGA 4 I 9.7). D'autre part, le foncteur $\xi$ est exact à gauche, comme $\varprojlim$ de foncteurs \inserted{repr. donc} exacts à gauche. Donc le fait que $\xi$ est représentable résulte du

\uline{Lemme} 4.8.11. Soit $M$ une catégorie accessible, \deleted{Alors les} \inserted{stable par $\varprojlim$ finies,} $M\xrightarrow{\xi}\operatorname{Ens}$ un foncteur. Pour que $\xi$ soit proreprésentable, il f. et s. qu'il soit exact à gauche et accessible. Donc \inserted{(SGA 4 I \unsure{8.12.2.1|8.12.8.1|§ 12.2.1})} si $M$ est stable par petites $\varprojlim$, $\xi$ est repr. \unsure{sss|ssi} il commute auxdites $\varprojlim$, et qu'il soit accessible.
% archive sheet number: 293
\scanpage{295}
\mspage{211}

Le ``il faut'' \unsure{est} \unsure{remis|mis} pour mémoire, on vient de le voir comme conséquence du fait que tout objet de $M$ est accessible. Le ``il suffit'' est essentiellement SGA 4 I 8.13.2. b), sauf que dans loc. cit on suppose $M$ \unsure{muni|munie} d'une \unsure{fc} \uline{stricte}. Mais la démonstration donnée dans loc. cit. marche aussi pour une \unsure{fcga}, en choisissant dans loc cit (p.137) $\pi$ tel que $\unsure{\xi|\zeta}$ soit $\pi$-accessible, et de plus $\pi$ soit \uline{\unsure{utile|stable}} pour la filtration cardinale gén. envisagée.

Reste à prouver $c''\overset{?}{\Longrightarrow}c$, sous réserve que ce soit bien vrai. Ce serait une conséquence du lemme suivant (où $D,\pi$ sont \unsure{remplacés|remplis} par $C,\pi_0$)

\uline{Lemme} 4.8.12 \unsure{(?)}. Soient $C$ une petite catégorie,

\margin{Lemme non prouvé \unsure{ici}, cf dém 4.11. (\unsure{en addenda|en admissib.})}% slanted

$\pi_0$ un cardinal infini. Si $C$ est stable par $\varprojlim$ finies, les $\operatorname{Ind}_{\pi_0}C$ également.

L'assertion similaire pour $\operatorname{Ind}(C)$, \unsure{ald|a.ld} $\operatorname{Ind}_{\pi}C$, est établie dans SGA 4 I 8.9.1 c). Mais il n'est pas clair sans plus que $\operatorname{Ind}_{\pi}C$ soit stable dans $\operatorname{Ind}(C)$ par $\varprojlim$ finies, dans le cas envisagé. \deleted{C'est lié à la} \unsure{D'en|Pour} \unsure{faire|bien} préciser, si $F=\varprojlim_{j\in J}F_j$, les $F_j$ dans $\operatorname{Ind}_{\pi}C$,
% archive sheet number: 294
\scanpage{296}
\mspage{212}

$F$ dans $\operatorname{Ind}(C)\subset C^\wedge$, on sait que $C_{/F}$ est filtrante, et il faut voir qu'\unsure{ille|il|elle} est filtrante grande devant $\pi_0$. On est ramené aux trois cas où $J=\varnothing$, $J=e\amalg e$, et $J=$

\diagram{
\begin{tikzcd}[ampersand replacement=\&]
\cdot\arrow[dr]\& \&\cdot\arrow[dl]\\
\&\cdot\&.
\end{tikzcd}
}

Les deux premiers cas sont essentiellement triviaux (car \unsure{$C$}, et un produit de deux cat. gdes devant $\pi_0$, sont gdes devant $\pi_0$), le troisième cas par contre \unsure{n'a|n'est} pas \inserted{l'air} \unsure{triviale|trivial} a priori. C'est lié visiblement à la question des représentations indicielles de foncteurs $J\longrightarrow\operatorname{Ind}_{\pi_0}(C)$, pour $\operatorname{card}\operatorname{red}J\leq\pi_0$. -- i.e. à la validité de 4.4.13 b), en dehors de l'hypothèse que $C$ soit \unsure{stable} par $\varinjlim$ de cardinal $\leq\pi_0$. Il nous faudra y revenir très prochainement, car cette question m'apparait à présent une question-clef, pour les propriétés de stabilité de la notion de catégorie accessible (en dehors d'\inserted{aucune} hypothèse d'existence de $\varinjlim$ ou $\varprojlim$).

\margin{\unsure{Z|2} \uline{Pas} \uline{prouvé} Mais \unsure{OK} si $\pi_0>\operatorname{card}\unsure{\operatorname{red}|\operatorname{sec}} C$ \unsure{$\Rightarrow$} cf. 4.15.10 p.~\unsure{463|462|403}.}% slanted

\uline{Proposition} 4.8.13. \deleted{\unsure{Soient}} Soit $C$ une catégorie équivalente à une petite catégorie, et \deleted{soit $\pi_0$} soit $\deleted{\unsure{C\geq}}\inserted{\pi_0\geq}\operatorname{card}\operatorname{red}C$ \deleted{un cardinal} \inserted{un cardinal} infini. Alors $M=\operatorname{Ind}_{\pi_0}C\deleted{\unsure{=}}=M_{\pi_0}$, \deleted{et c'est} (i.e. tous les objets de $\operatorname{Ind}_{\pi_0}C$ sont $\pi_0$-accessibles) et $M$ est une enveloppe de Karoubi de $\unsure{\widetilde C|\overline C|C}$,
% archive sheet number: 295
\scanpage{297}
\mspage{213}

Signalons tout de suite les corollaires suivants.

\uline{Corollaire} 4.8.14. Soit $M$ une $U$-catégorie. Conditions équivalentes :

\margin{\unsure{Z|2} \uline{Pas} \uuline{prouvé} Si \unsure{finalemt}, cf 4.15.10 et 4.15.11. p.~462/463}% slanted

a) $M$ est accessible, et est équivalente à une petite catégorie.

$a')$ $M$ est accessible, et il existe un cardinal infini $\pi_0$ tel que tous les objets de $M$ soient $\pi_0$-accessible\unsure{s} \deleted{[De plus, si $\pi_0$ est \unsure{fi}} \inserted{(i.e. $M=M_{\pi_0}$)}

\deleted{$a'')$ \unsure{(b}} $b)$ $M$ est équivalente à une petite catégorie, et est Karoubienne i.e. stable par facteurs directs.

\margin{NB Ceci implique que $M$ et $M^\circ$ sont accessibles, et je conjecture que la réciproque est vraie également.}

L'équivalence de $a)$ et $a')$ est \unsure{clair|claire}, puisque $M$ admet une filtration cardinale \inserted{gén. canonique} $(M_{\pi})_{\pi\geq\pi_0}$, formée de \deleted{petites} catégories équivalentes à des petites catégories \inserted{(condition \unsure{a)} des \unsure{fcg})}, et telle que toute telle sous-catégorie soit contenue dans une des $M_{\pi}$ (\inserted{résulte de la} condition b) des fcg). Comme $M_{\pi_0}$ est équivalente à une petite catégorie et est karoubienne (étant stable par $\varinjlim$ filtrantes grandes devant $\pi_0$), il est clair aussi que $a')\Longrightarrow b)$. Reste à prouver $b)\Longrightarrow a')$, or \deleted{si $\pi_0$} cela résulte aussitôt de 4.8.13 \inserted{appliqué à $M$}, qui montre de plus qu'on peut prendre pour $\pi_0$ n'importe quel cardinal infini tel que $\pi_0\geq\operatorname{card}\operatorname{red}\deleted{\unsure{C}} M$.
% archive sheet number: 296
\scanpage{298}
\mspage{214}

\inserted{4.8.15.} \uline{Dém. de 4.8.13}. Nous allons \deleted{\unsure{d'abord}} prouver :

\uline{Lemme 4.8.15.1}. Sous les conditions de 4.8.13, soit $F\in\operatorname{Ob}M$, $\deleted{\unsure{\text{alors}|x\in\operatorname{Ob}}}\inserted{x\in\operatorname{Ob}}C$, alors

\[
\operatorname{card}\deleted{\unsure{\operatorname{card}|x}}F(x)\leq\pi_0.\tag*{($\ast$)}
\]

Ceci signifie donc, en vertu de 4.7.1 $a''\Longrightarrow b$ pour $\pi\geq\pi_1$, compte tenu qu'ici $\pi_{\deleted{\unsure{0}}\unsure{1|01}}=\pi_0$, $\pi=\pi_0$, qu'on a $F\in\operatorname{Filt}^{\pi_0}$ d'où a fortiori $F\in M_{\pi_0}$, ce qui \deleted{prouve\unsure{ra}} implique $M=M_{\pi_0}$ i.e. la proposition 4.8.13, \inserted{compte tenu de 4.8.5.} Donc tout revient à prouver $(\ast)$.

Écrivons

\[
F=\varinjlim_{i\in I}x_i\qquad \left|\begin{aligned}&\text{les }x_i\in\operatorname{Ob}C\\ &I\ \text{ens. ordonné}\\ &\text{grd devant }\pi_0,\end{aligned}\right.
\]

d'où

\[
F(x)=\varinjlim_{i\in I}\operatorname{Hom}_C(x,x_i).
\]

OPS

\[
\operatorname{card}C\leq\pi_0,
\]

quitte à remplacer $C$ par une catégorie réduite équivalente.

Pour $\forall\,\unsure{\xi|\zeta}\in\operatorname{Ob}C$, soit \deleted{$\unsure{I(\xi)}$}

\[
I(\xi)=\{i\in I\mid x_i=\xi\},
\]

\[
\text{alors}\qquad I=\bigcup_{\xi\in\operatorname{Ob}C}I_\xi,\qquad\deleted{\unsure{\text{et}}}
\]
% archive sheet number: 297
\scanpage{299}
\mspage{215}

réunion d'une famille de parties \inserted{de $I$} de cardinal $\operatorname{card}\operatorname{Ob}C\leq\pi_0$. Comme $I$ est grd devant $\pi_0$, on voit que cela implique que l'une des $I_{\unsure{\xi|\zeta}}$ est \unsure{cofinal|cofinale} dans $I$. Quitte à remplacer $I$ par $I_\xi$, OPS $I_\xi=I$, i.e. que tous les $x_i$ sont égaux à un \unsure{même} élément $\xi$. Quitte à remplacer $C$ par la sous-catégorie pleine $\deleted{\unsure{\xi}}\inserted{C_0=} \{\xi\}$, OPS $C$ réduite à un seul élément. \unsure{[|(}Car si on prouve que $F$, en tant qu'objet de $\operatorname{Ind}_{\pi_0}(C_0)$, est facteur direct d'un objet de $C_0$, (i.e. de $\xi$), il en sera de même de $F$ en tant qu'objet de $\operatorname{Ind} C$). Donc OPS

\[
C=\{\xi\},
\]

donc $C$ donné par le monoïde

\[
M=\operatorname{End}_C(\xi),\qquad \operatorname{card} M\leq\pi,
\]

et la donnée du syst. inductif $(x_i)_{i\in I}$ revient à la donnée d'un système d'éléments

\[
\alpha_{ji}\in M\qquad (i,j\in I,\ i\leq j)
\]

\[
\text{satisfaisant}\qquad\alpha_{ii}=1,\quad\alpha_{hj}\alpha_{ji}=\alpha_{hi}\quad\text{si }i\leq j\leq h.
\]

On a

\[
F(\xi)=\varinjlim_{i\in I}M_i\qquad \left|\begin{aligned}&\text{où }M_i=M\quad\forall\,i\in I,\\&\text{et }M_i\longrightarrow M_j\ \text{est}\\&\text{donné par}\\ &\qquad u\longmapsto\alpha_{ji}u\end{aligned}\right.
\]
% archive sheet number: 298
\scanpage{300}
\mspage{216}

\margin{Pour ces réductions, cf 4.15.8 (p.453 à \unsure{468|463}). L'ultime réduction se trouve dans \unsure{4.15.8.9} (p.~\unsure{460|468}), et son corollaire 4.15.10.}% slanted

J'ai pu me ramener, par des arguments similaires,$\vphantom{x}^{\ast}$ au cas où le monoïde $M=\operatorname{End}\unsure{(\xi)|(\zeta)}$ est un monoïde de monomorphismes i.e.

\[
uv=uv'\Longrightarrow v=v'\qquad\text{si }u,v,v'\in M,
\]

ce \deleted{\illegible{1}} qui implique que les flèches de transition dans le système inductif des $M_{\unsure{i|ki}}=M$ sont des mono. Il s'ensuit que l'on ne peut avoir $\operatorname{card}F(\xi)\leq\pi_0$ ($I$ étant grand devant $\pi_0$) que si $\alpha_{j i_0}$ est un iso pour \inserted{$j\geq i_0$, et} $i_0$ assez grand, i.e. si le syst. inductif \deleted{est \unsure{ess. co}} induit \deleted{par} pour $\deleted{\unsure{i}}j\geq i_0$ est constant, a fortiori que le syst. inductif envisagé est \unsure{ess.t|esst.} constant. Mais je viens de passer près de quatre heures pour essayer de prouver ce point, sans y être arrivé ! Cela semble plus délicat à prouver que je ne le prévoyais, mais je crois toujours que c'est bel et bien vrai.

\margin{$\unsure{\xi|\rbrace}$}
% archive sheet number: 299
\scanpage{301}
\mspage{217}

4.9. \uline{Le théorème de \unsure{représentations indicielles|représentation indicielle}} \uline{pour} $\underline{\operatorname{Ind}}_{\pi}(C)$.

4.9.1. \unsure{Soient} $C$ une catégorie équivalente à une petite catégorie, $\pi_0$ un cardinal infini,

\[
M=\operatorname{Ind}_{\pi_0}(C)\subset C^\wedge,\tag*{(4.9.1)}
\]

$J$ une catégorie équivalente à une petite catégorie. Considérons l'inclusion \inserted{pl. fid} $C\hookrightarrow M$, et l'inclusion correspondante de $U$-catégories

\[
\underline{\operatorname{Hom}}(J,C)\hookrightarrow\underline{\operatorname{Hom}}(J,M).\quad\text{(pl. fid.)}.\tag*{(4.9.2)}
\]

On sait, \deleted{que} comme $C\subset M_{\pi_0}$, que les objets de $\underline{\operatorname{Hom}}(J,C)$ donnent des objets $\pi_0$-\uline{accessibles} de \deleted{$\unsure{C}$,} $\underline{\operatorname{Hom}}(J,M)$, du moins si on suppose

\[
\operatorname{card}\operatorname{red}J\leq\pi_0,\tag*{(4.9.\unsure{3|B})}
\]

ce que nous ferons désormais. (Cf 4.5.8, \unsure{1°|?°}, p.~119) Comme le deuxième \unsure{mb} de (4.9.2) satisfait $L_{\pi_0}$, ($M$ y satisfaisant), il s'ensuit que \deleted{l'in} le foncteur (4.9.2) se prolonge en

\[
\operatorname{Ind}_{\pi_0}(\underline{\operatorname{Hom}}(J,C))\hookrightarrow\underline{\operatorname{Hom}}(J,\underset{\substack{\Vert\\ \operatorname{Ind}_{\pi_0}(C)}}{M}),\tag*{(4.9.3)}
\]

et ce foncteur est \uline{pleinement fidèle} grâce au résultat d'accessibilité qu'on vient d'invoquer. Je m'intéresse à établir des cas où ce foncteur est essentiellement surjectif, i.e. est une \uline{équivalence} de catégories. Notons que pour que \inserted{$\varphi:J\unsure{\longrightarrow|\text{---}}M$} appartienne à l'image essentielle, il f. et s. qu'il existe un foncteur

\[
\Phi:I\longrightarrow\underline{\operatorname{Hom}}(J,C)\qquad\unsure{\Biggl\lbrace|\Biggl\vert}\begin{aligned}&I\ \text{grd devant }\pi_0\\ &\text{(et OPS }I\ \\&\text{ordonné)}\end{aligned},\tag*{(4.9.4)}
\]
% archive sheet number: 300
\scanpage{302}
\mspage{218}

tel qu'on ait

\[
\varphi\xleftarrow{\sim}\varinjlim_{i\in I}\Phi(i),\tag*{(4.9.5)}
\]

où ce qui revient au \unsure{même}, tel qu'on ait

\[
\varphi(j)\xleftarrow{\sim}\varinjlim_{i\in I}\Phi(i,j),\quad\text{pour }j\in J\tag*{(4.9.6)}
\]

où on a posé

\[
\Phi(i,j)=\Phi(i)(j)
\]

interprétant donc $\Phi$ comme un foncteur

\[
\Phi:I\times J\longrightarrow C\ ;\tag*{(4.9.7)}
\]

la loi fonctorielle de $\varphi(j)$ donnée par (4.9.6) \unsure{étant|est} évidente. En d'autres termes, les $\varphi$ dans l'image essentielle de (4.9.3) sont ceux qui admettent une ``\uline{représentation indicielle}'' par un foncteur $\Phi$ sous forme (4.9.4), ou sous forme (4.9.7), avec \uline{de plus} $I$ grd devant $\pi_{\unsure{0|1}}$. Je rappelle à ce sujet la

\uline{Proposition} 4.9.2. Sous les conditions précédentes sur $C,\pi_0,J$ (notamment $\operatorname{card}\operatorname{réd}J\leq\pi_0$), si $\varphi:J\longrightarrow M\overset{\mathrm{\unsure{def|déf}}}{=}\operatorname{Ind}_{\pi_0}(C)$, pour que $\varphi$ appartienne à l'image essentielle du foncteur pl. fid. (4.9.3), il f. et s. qu'elle satisfasse les deux conditions suivantes :
% archive sheet number: 301
\scanpage{303}
\mspage{219}

a) $\underline{\operatorname{Hom}}(J,C)_{/\varphi}$ est grde devant $\pi_0$.

b) Pour tout $j\in J$, le foncteur

\[
\underline{\operatorname{Hom}}(J,C)_{/\varphi}\longrightarrow C_{/\varphi(j)}\qquad\left|\begin{aligned}&\text{foncteur entre}\\&\text{catégories filtrantes}\\&\text{grdes devant }\pi_0\\&\text{en vertu de a)}\end{aligned}\right.\tag*{(4.9.2.1)}
\]

induit par le foncteur

\[
\left|\begin{aligned}\underline{\operatorname{Hom}}(J,C)&\longrightarrow C\\ F&\longmapsto F(j)\end{aligned}\right.,\tag*{(4.9.2.2)}
\]

est \uline{cofinal}.

Le résultat principal \deleted{du}\inserted{de} \unsure{la} présente section est le suivant :

\uline{Théorème} 4.9.3. Sous les conditions générales de 4.9.1 pour $C,\pi_0,J$ (notamment $\operatorname{card}\operatorname{réd}J\leq\pi_0$), le foncteur pl fidèle (4.9.3) est une équivalence de catégories dans chacun des deux cas suivants :

a) $C$ est stable par $\varinjlim$ filtrantes dénombrables.

\margin{Cas \unsure{commun} : il existe un ordinal limite $\alpha$, tel que 1) \deleted{Pour \unsure{tt entier} $\beta<\alpha$, \illegible{1}} $\operatorname{card}\deleted{\unsure{\beta}}\alpha\leq\unsure{\pi|\pi_0}$. 2) les limites inductives de type $I_\alpha=\{\beta\mid\beta<\alpha\}$ \unsure{st|sont} représentables dans $C$}% slanted

b) Il existe un cardinal $\pi'_0<\pi_0$, et un cardinal $\pi''_0>\pi'_0$ (NB OPS que $\pi''_0$ est le successeur de $\pi'_0$)\inserted{donc $\pi_{\!0}^{\unsure{'|''}}\leq\pi$}, tels que $C$ soit stable par $\deleted{\varinjlim}$limites inductives \unsure{grdes} devant $\pi'_0$, et de cardinal $\leq\pi''_0$.

\uline{Remarques} 4.9.4. 1) Ce théorème, cas a), contient le th. 4.4.13, partie b) (p.~95-96), où on suppose que $C$ est stable par limites inductives de cardinal $\leq\pi_0$, l'hypothèse a) de \unsure{4.9.3} étant considérablement plus faible.
% archive sheet number: 302
\scanpage{304}
\mspage{220}

b) Je présume qu'en dehors d'une hypothèse restrictive, du type a) ou b) ci-dessus, le foncteur \inserted{pl fid.} (4.9.3) n'est pas tjrs ess. surjectif, même dans le cas particulier où $J$ est une catégorie finie, p.ex. $J=B_G$, $G=\{\pm1\}$. Pour un contre-exemple dans le cas très voisin où on prend $M=\operatorname{Ind}(C)$ (au lieu de $\operatorname{Ind}_{\pi_0}(C)$), et où on considère le foncteur \inserted{(pl. fid. si $J$ est une catégorie équiv. à une cat. \uline{finie})}

\[
\operatorname{Ind}(\underline{\operatorname{Hom}}(J,C))\longrightarrow\underline{\operatorname{Hom}}(J,\operatorname{Ind}(C)),
\]

je renvoie à l'exercice SGA 4 I 8.8.8.

\unsure{3)|c)} L'application principale du th. 4.9.3. consiste en le

\uline{Corollaire} 4.9.5. Soient $M$ une catégorie \uline{accessible} (4.8.7), $J$ une catégorie équivalente à une petite catégorie. Alors la catégorie \inserted{$M(J^\circ)=$} $\underline{\operatorname{Hom}}(J,M)$ est accessible. Plus précisément, soient $\pi_0$ un cardinal \uline{utile} pour $M$ (i.e tel que $\operatorname{Ind}_{\pi_0}(M_{\pi_0})\xrightarrow{\approx} M$, \deleted{\unsure{) et}} \deleted{soit $\pi$} i.e. \unsure{tout} cardinal tel que $M$ soit $\pi_0$-accessible), soit $\pi$ un cardinal \inserted{infini} \deleted{\unsure{tel}} que

\[
\pi^{\pi_0}=\pi\geq\operatorname{card}\operatorname{red}J,\quad\overbrace{\deleted{\text{(et }\pi\unsure{\geq|>}\operatorname{card}\operatorname{red} M_{\pi_0}}}^{\inserted{\ast}}\quad\deleted{\text{p. ex.}}\inserted{\text{i.e.}}\quad\pi=\deleted{\unsure{C}}c^{\pi_0},\quad\text{où }\deleted{\unsure{\pi}}
\]

\margin{(On a donc $\pi\deleted{\unsure{\geq}}=\pi^{\pi_0}\geq 2^{\pi_0}>\pi_0$)}% slanted

\deleted{$\pi'_{\unsure{1|g}}=\operatorname{Sup}(\pi_0,\operatorname{card}\operatorname{red}J)$ \unsure{est tel que}} $c$ est n'importe quel cardinal $c\geq\operatorname{Sup}(\pi_0,\operatorname{card}\operatorname{red}J)$, alors $\underline{\operatorname{Hom}}(J,M)$ est $\pi$-accessible, plus précisément on a une équivalence de catégories

\[
\underline{\operatorname{Hom}}(J,M)\xleftarrow{\approx}\operatorname{Ind}_{\pi}(\underline{\operatorname{Hom}}(J,M_{\pi}))\tag*{(4.9.5.1)}
\]
% archive sheet number: 303
\scanpage{305}
\mspage{221}

On sait en effet, \unsure{sous|avec} les hypothèses faites sur $\pi_0,\pi$, que \deleted{\unsure{l'on a}} $\pi$ est utile pour $M$ i.e que

\[
M\xleftarrow{\approx}\operatorname{Ind}_{\pi}(M_{\pi}),
\]

d'autre part $M_\pi$ est stable par $\varinjlim$ grdes devant $\pi_0$ et de cardinal $\leq\pi$. On peut donc appliquer 4.9.3 b), en prenant $C=M_{\pi}$, $\pi$ et $\pi_0$ comme ici (donc $\pi_0<2^{\pi_0}\leq\pi$), et $\pi'_0=\pi$.

\margin{Résultats préliminaires (à la dém. de \unsure{4.9.3|4.9.7})}% slanted

4.9.6. Le reste de la présente section 4.9. est \unsure{consacrée|consacré} à la démonstration de 4.9.3, i.e. à la vérification, \deleted{pour} $\varphi:J\longrightarrow M$, des conditions a), b) de 4.9.2. Dans toute la suite, on se fixe donc \unsure{un|une}

\[
\varphi:J\longrightarrow M=\operatorname{Ind}_{\pi}(C).\tag*{(4.9.\unsure{6|5}.1)}
\]

\inserted{On peut supposer $J$ réduite, donc petite et \uline{$\operatorname{card}J\leq\pi$}}% slanted

On va considérer le foncteur déduit des foncteurs (4.9.2.1)

\[
\underline{\operatorname{Hom}}(J,C)_{/\varphi}\longrightarrow\prod_{j\in\operatorname{Ob}J} C_{/\varphi(j)},\tag*{(4.9.6.2)}
\]

dont les foncteurs en question se \unsure{déduisant|déduisent} en composant avec les foncteurs projection

\[
\prod C_{/\varphi(j)}\longrightarrow C_{/\varphi(j_0)}.\tag*{(4.9.6.3)}
\]

Comme les catégories facteurs $C_{/\varphi(j)}$ sont filtrantes (et \unsure{même} grdes devant $\pi$), leur produit est également filtrant, et même grd devant $\pi$ (immédiat), et de même
% archive sheet number: 304
\scanpage{306}
\mspage{222}

pour les produits partiels. Il s'ensuit aussitôt que les projections (4.9.6.3) sont cofinales, en vertu du

\uline{Lemme} 4.9.6.1. a) Soient $I,I'$ deux catégories$\vphantom{x}^{\unsure{\ast}}$ \deleted{filtrantes \unsure{avec $I$} filtrantes,} \inserted{avec $\unsure{I'|I}$ 0-connexe (p.ex. \unsure{filtrant|filtrante}),} alors le foncteur projection $\operatorname{pr}_2:I\times I'\longrightarrow I'$ est cofinal.

\deleted{b) Tout produit de cat. filtrantes est \inserted{cat.} filtrante, (donc 0-connexe)}

\deleted{\unsure{a)}} En effet, pour $i'\in\operatorname{Ob}I'$, on a

\[
I\times I'\simeq {}_{\unsure{i'\backslash|i\backslash}}I\times{}_{\unsure{i'\backslash|i'/}}I',
\]

$\vphantom{I}_{i'\backslash}I$ et $\vphantom{I}_{i'\backslash}I'$ étant 0-connexes, il en est de \unsure{même} de leur produit.

\deleted{b) Vérification immédiate, en \unsure{PF1, PF2}.}

\margin{Attention, il n'est \uline{pas} vrai qu'un produit infini de cat.\textsuperscript{g} 0-connexes \unsure{soit|sont|est} 0-connexe ! $\unsure{(\ast)}$}% slanted

\uline{\deleted{\unsure{Corollaire|Lemme}}\unsure{Lemme|Corollaire}} 4.9.6.2. Soit $(I_\alpha)_{\alpha\in A}$ une famille de catégories \uline{filtrantes} \inserted{(resp. grdes devant $\pi$)}. Alors \deleted{la \unsure{projection}} $I_A\overset{\mathrm{def}}{=}\prod_{\alpha\in A}I_\alpha$ est filtrante (resp gde devant $\pi$), et pour $\forall$ partie $B\subset A$, la projection canonique $\deleted{\unsure{I_A\longrightarrow\prod_{\alpha\in B}I_\alpha}}$ $I_A\longrightarrow I_B$ est cofinale.

\margin{Cf dém élégante ci-dessous, 4.9.6.\unsure{6|5}.}% slanted

La stabilité de la propriété ``filtrante'' (ou ``gde devant $\pi$'' par \unsure{AQT|AQI|$\forall\prod$}. On a, si $B'=A\setminus B$, $\quad I_A\simeq I_B\times I_{B'}\longrightarrow I_B$, et on applique le lemme 4.9.6\unsure{,|.}1 pour conclure que ce foncteur est cofinal, $I_{B'}$ étant filtrante donc 0-connexe.
% archive sheet number: 305
\scanpage{307}
\mspage{223}

Le foncteur (4.9.6.3) étant cofinal, pour établir que \unsure{4.9.2.1)} l'est, il suffit de prouver qu'il en est ainsi du foncteur (4.9.6.2), \inserted{ce que nous ferons.} Notons de plus que ledit foncteur est \uline{fidèle}, car il s'insère dans le carré commutatif

\diagram{
\begin{tikzcd}[ampersand replacement=\&]
\underline{\operatorname{Hom}}(J,C)_{/\varphi} \arrow[r] \arrow[d] \& \prod C_{/\varphi(j)} \arrow[d] \\
\underline{\operatorname{Hom}}(J,C) \arrow[r] \& \deleted{\unsure{\prod|F}} C^{\operatorname{Ob}J}
\end{tikzcd}
}

où les flèches verticales et la deuxième flèche horizontale sont fidèles. \deleted{\unsure{On va|En ce}} \deleted{alors} En définitive, pour établir le th. \unsure{4.9.3}, nous allons étudier le foncteur \inserted{\uline{fidèle}} (4.9.6.2)

\[
f\colon\underline{\operatorname{Hom}}(J,C)_{/\varphi}\unsure{\longrightarrow|\text{---}}\prod_{j\in J_0}C_{/\varphi(j)}\tag*{($\ast$)}
\]

(où on a posé $J_0=\operatorname{Ob}J$, $J_1=\mathrm{Fl}\,J$), et \deleted{il \unsure{faut|faudra} nous} nous allons prouver :

\textcircled{a} \deleted{\unsure{Le}}Ce foncteur \inserted{fidèle} $(\ast)$ est \uline{cofinal}.

\textcircled{b} La catégorie-source \inserted{$\underline{\operatorname{Hom}}(J,C)_{/\varphi}$} est gde devant $\pi_0$.

Nous allons pour ceci dégager un lemme \unsure{auxiliaire|critérium|condition} sur la cofinalité des foncteurs \uline{fidèles} :
% archive sheet number: 306
\scanpage{308}
\mspage{224}

\margin{Ce lemme fait partie des \unsure{réslts|résultats} généraux sur la cofinalité et les cat. filtrantes.}% slanted

\uline{\deleted{\unsure{Proposition}}} 4.9.6.3. Soit $f\colon I^{\unsure{a|4}}\unsure{\text{---}|\longrightarrow}I'$ un foncteur, et considérons les conditions :

a) $I$ filtrante, $f$ cofinal i.e. \inserted{(en présence de $I$ \unsure{filtr.te|filtr.-e})} $f$ satisfait les conditions fonctorielles $\text{\textcircled{F1}},\text{\textcircled{F2}}$.

\inserted{b) $I'$ est filtrante, et les $\vphantom{I}_{i'}\backslash I$ (pour $i'\in\operatorname{Ob}I'$) sont filtrantes.}

\unsure{c)|b)} $I'$ filtrante, et $f$ satisfait à F1, ainsi qu'à la condition suivante

\textcircled{F3} Pour tout couple $(i_0,i_1)$ d'objets de $I$ et toute flèche \deleted{\unsure{$u$:}} $f(i_0)\xrightarrow{u'}f(i_1)$ il existe un objet $i_2$ de $I$, et une flèche $v'\colon f(i_1)\longrightarrow f(i_2)$, tels que \deleted{\unsure{les|des}} $v'$ et $v'u'$ \deleted{\unsure{puissent s'écrire}} \deleted{\unsure{proviennent des}} \inserted{soient relevables} \unsure{en|des} flèches \deleted{\unsure{de}} $i_{\unsure{1|0}}\xrightarrow{v}i_2$ et $i_0\overset{w}{\unsure{\text{---}|\longrightarrow}}i_2$ \unsure{;|:}

\diagram{
\begin{tikzcd}[ampersand replacement=\&]
\unsure{i_0|j_0} \arrow[r,"{\deleted{\unsure{u}}}",dotted] \arrow[rr,"w",bend left=21,dotted] \arrow[d,no head] \& i_1 \arrow[r,"v",dotted] \& i_2 \\
f(i_0) \arrow[r,"{u'}"] \arrow[rr,"{f(w)}"',bend right=22,dashed] \& f(i_1) \arrow[r,"{f(v)=v'}",no head] \& f(i_2)
\end{tikzcd}
}

$[$ i.e il existe $\unsure{i_{2a}|i_2}$ et $v\colon i_1\unsure{\longrightarrow|\text{---}}i_2$, $w\colon i_0\unsure{\longrightarrow|\text{---}}i_2$ tels qu'on ait $f(v)u'=f(w)$ $]$.

\deleted{\unsure{c')} $I'$ filtrante, les \unsure{$I_{i'}$} et les $\vphantom{I}_{i'}\backslash I$ \unsure{ps.|pseudo-}filt.}

\inserted{b') $I'$ $0$-connexe, les $\vphantom{I}_{i'}\backslash I$ filtrantes.}

\inserted{b'') $I'$ est filtrante, \inserted{satisfait PF$(2)^*$ (égalisation des doubles flèches)} les $\vphantom{I}_{i'}\backslash I$ ps. filtr\unsure{antes|.}, \deleted{\unsure{les}} \unsure{i.e.} satisfont à la condition PF$1'$.}

\margin{p. ex. PF\unsure{3|2} est satisfaite si $I'$ est \uline{ordonnée}, sans avoir à supposer alors de $I'$ filtrante, sauf \deleted{\unsure{ment}} \uline{$0$-connexe}}% slanted

Ceci posé :

1) On a les implications

\[
\begin{aligned}&\deleted{\unsure{a\Longleftrightarrow b|a\Longleftarrow b}}\unsure{\Longrightarrow|\Longleftrightarrow}c\\
&a\Longleftrightarrow b\Longleftrightarrow b'\Longrightarrow b''\unsure{\Longleftrightarrow|\Longleftarrow}c\quad\deleted{\unsure{b\Longleftrightarrow c|\Longleftrightarrow a}}\end{aligned}
\]

et si $f$ est \uline{fidèle} \inserted{(pour $I$ \unsure{ordonnée|ndonnée})}, les \deleted{trois} \inserted{\unsure{cinq|quatre}} conditions a), b), \inserted{b'), b''),} c) sont équivalentes.

\deleted{2) \unsure{plus|pour} \unsure{$\pi$|R} \unsure{un|un card.} \unsure{cardinal infini :} \unsure{et} condition : \unsure{équiv|équivalente|Supposons}} \deleted{\unsure{Pour que} Supposons $I$ grande devant $\pi$,} \deleted{et $f$ cofinal, alors $I'$ est grande devant $\pi$}

$(a_\pi)$ $I$ gde \unsure{dev.t|dev.} $\pi$, $f$ cofinal

$(b_\pi)$ $I'$ grande devant $\pi$, les $\vphantom{I}_{i'}\backslash I$ aussi.

Alors $a_\pi\Rightarrow b_\pi$, et si $f$ fidèle, $a_\pi$ et $b_\pi$ st équiv.
% archive sheet number: 307
